\section{Discussion}
\label{sec:discussion}

The results in Section~\ref{sec:results} support a narrower claim than the architecture's original
design intuition promised, and the narrower claim is the more useful one. ``The rule engine catches
the model's mistakes'' turns out to describe only 12\% of what actually happens; ``the system
defaults to a safe state whenever the model's output cannot be trusted, and never lets a trusted
output fall below what the rule engine alone would say'' describes all of it. That the first,
catchier framing is wrong is not a weakness in the finding -- it is the finding. A reader who
walks away believing the rule engine spends most of its time correcting nuanced severity
misjudgments would misunderstand what this system's safety property actually rests on, and would
likely misjudge how well the same pattern would hold up in their own deployment if they assumed
the model would usually produce something parseable to correct in the first place.

\subsection{What a better-tuned model would and would not change}

The most obvious objection to the central result is that AQUA-1B was tested under close to the
worst plausible version of itself: an 8-example adapter its own training pipeline calls
insufficient, with a properly sized alternative sitting unused. Whether the documented 128-example
adapter would reduce the 50\% unparseable rate is genuinely untested, and this paper does not
claim to know the answer -- it is stated here as open future work, not as a hedge against the
current result. What is not open, and would not change regardless of which adapter is loaded, is
the architecture's independence from the model's free-text reasoning: \texttt{apply\_rule\_override}
never reads the \texttt{reasoning} field under either adapter, so a better-tuned model could raise
the agreement rate without touching the property this paper actually measured. The two questions
-- does the model produce more parseable, more accurate output, and does the safety property hold
regardless -- are separable by design, and conflating them would understate how load-bearing the
narrow-interface choice actually is.

A second, related question this system's own agentic-routing result raises but cannot answer:
whether the same escalation-only, narrow-interface discipline holds for agentic tool-routing at a
larger on-device model scale. Both models tested here either failed outright (AQUA-1B) or produced
a constant-answer artifact (AQUA-7B); neither result licenses a claim about where on the capability
curve a model would need to sit before deterministic routing stops being the safer choice. That
threshold, if it exists, was not found here, because it was not this paper's question to answer --
the negative result stands on its own terms, for these two models, on this scenario set.

\subsection{The sparse-frame gap, stated as a gap}

The reassuring empirical result that zero of 98 validation frames produced a full-frame miss is
conditional on a property of the validation set that nobody designed in: every one of those 98
frames contains at least three ground-truth fish. Whether the same detector would miss a genuinely
sparse frame -- one or two fish, the case a welfare check most needs to catch, since an isolated or
missing fish is itself a potential welfare signal -- is a question this dataset cannot answer in
either direction. The extrapolated 22--28\% miss-risk estimate at $k=1$ is exactly that: an
extrapolation from an independence approximation, not a measurement, and it should not be read
with the same confidence as the empirical 0/98 result it sits beside. Closing this gap requires
collecting or synthesizing frames the current dataset simply does not contain, and that is stated
here as a concrete, actionable follow-up rather than a general caveat about validation-set size.

\subsection{Practitioner rules this paper's evidence actually supports}

Four rules survive scrutiny against the specific evidence behind them, and are stated here at that
scale rather than inflated into general design advice.

\emph{Architect a safety-critical LLM-plus-validator pipeline assuming the model will frequently
produce nothing usable, verify that your fallback default is the conservative one, and never let
your validator read anything from the model beyond the one enumerable field it is designed to
certify.} This is the paper's central practitioner takeaway, and it follows directly from
Section~\ref{sec:results-dual-layer}: the property that held was never ``the model reasons well
enough,'' it was ``the validator only ever looks at one field, and that field defaults safe.''

\emph{Tune a retrieval threshold for hallucination-averse domains against asymmetric cost, not raw
F1.} Section~\ref{sec:results-rag}'s 0.85-over-0.84 choice sacrifices measurable recall
specifically to drive hard-negative pass-through to zero, because a fabricated, confidently cited
context actively misinforms while a missed document only weakens reasoning a downstream
deterministic layer still adjudicates. The transferable lesson is the cost function a practitioner
should write down before tuning, not this specific threshold value.

\emph{Choose a drift-detection binning scheme by sweeping your own confidence distribution's
severity range, not by copying a convention.} Section~\ref{sec:results-mlops}'s crossover means
``use quantile binning'' is not a universal rule -- it is a consequence of this particular
detector's confidence outputs being concentrated in a narrow band. A practitioner adopting PSI for
a differently shaped output distribution should check for the same crossover before assuming either
binning convention behaves as expected across their own severity range.

\emph{Benchmark the whole exported inference pipeline before attributing an accuracy loss to
quantization.} Section~\ref{sec:results-vision}'s INT8-loses-less-than-fp32 finding means the
standard mental model -- lower numeric precision costs more accuracy -- can be wrong for a given
toolchain when the real loss sits in a shared post-processing path common to every export target.
The actionable version of this rule is to measure the full configuration matrix, not to reason
about precision loss in isolation.

A fifth pattern is worth naming even though it is not a design rule so much as a discipline: the
adapter-provenance finding and the live 0.695 leak inside the system's own RAG knowledge base
(Section~\ref{sec:results-audit}) are the same phenomenon at two different layers of the stack --
an artifact silently drifted, or was mis-identified, between being produced and being used to
support a conclusion. Both were caught by cheap, simple file-level checks (a record count and an
mtime comparison; a repository-wide grep) that most teams skip because ``we deployed the right
thing'' is treated as an assumption rather than a fact worth checking.

\subsection{A hedged conjecture}

These results suggest, but do not demonstrate beyond this one system, that narrow-enumerable-field
validator interfaces may be a generally transferable safety-design principle for LLM-plus-validator
pipelines outside aquaculture monitoring. The support for this is the consistency of the same
underlying discipline across five architecturally distinct layers measured here -- decision fusion,
agentic tool routing, retrieval, MLOps retraining, and edge vision -- not any test conducted outside
this one system. This is offered as a direction worth testing in other domains, not as a
demonstrated generalization, and a reader should weight it accordingly.
